\documentclass[10pt,a4paper]{amsart}
\usepackage{amsmath,amssymb,amsthm,graphicx}
% Retained minimal preamble for the Stacks Project chapter varieties.tex.
% Verbatim excerpt of the chapter preamble: Xy-pic setup (native preamble lines 14-18), the native theorem
% declarations with the mathematical macros (native preamble lines 154-217), and the project own no-op
% definitions of the reference/slogan/history display environments (native wrapper.tex lines 1-3) that the
% counted statement uses for its slogan. Copied mechanically; nothing added, removed or redefined.
% For commutative diagrams we use Xy-pic
\usepackage[all]{xy}

% We use 2cell for 2-commutative diagrams.
\xyoption{2cell}
% Theorem environments.
%
\theoremstyle{plain}
\newtheorem{theorem}[subsection]{Theorem}
\newtheorem{proposition}[subsection]{Proposition}
\newtheorem{lemma}[subsection]{Lemma}

\theoremstyle{definition}
\newtheorem{definition}[subsection]{Definition}
\newtheorem{example}[subsection]{Example}
\newtheorem{exercise}[subsection]{Exercise}
\newtheorem{situation}[subsection]{Situation}

\theoremstyle{remark}
\newtheorem{remark}[subsection]{Remark}
\newtheorem{remarks}[subsection]{Remarks}

\numberwithin{equation}{subsection}

% Macros
%
\def\lim{\mathop{\mathrm{lim}}\nolimits}
\def\colim{\mathop{\mathrm{colim}}\nolimits}
\def\Spec{\mathop{\mathrm{Spec}}}
\def\Hom{\mathop{\mathrm{Hom}}\nolimits}
\def\Ext{\mathop{\mathrm{Ext}}\nolimits}
\def\SheafHom{\mathop{\mathcal{H}\!\mathit{om}}\nolimits}
\def\SheafExt{\mathop{\mathcal{E}\!\mathit{xt}}\nolimits}
\def\Sch{\mathit{Sch}}
\def\Mor{\mathop{\mathrm{Mor}}\nolimits}
\def\Ob{\mathop{\mathrm{Ob}}\nolimits}
\def\Sh{\mathop{\mathit{Sh}}\nolimits}
\def\NL{\mathop{N\!L}\nolimits}
\def\CH{\mathop{\mathrm{CH}}\nolimits}
\def\proetale{{pro\text{-}\acute{e}tale}}
\def\etale{{\acute{e}tale}}
\def\QCoh{\mathit{QCoh}}
\def\Ker{\mathop{\mathrm{Ker}}}
\def\Im{\mathop{\mathrm{Im}}}
\def\Coker{\mathop{\mathrm{Coker}}}
\def\Coim{\mathop{\mathrm{Coim}}}

% Boxtimes
%
\DeclareMathSymbol{\boxtimes}{\mathbin}{AMSa}{"02}

%
% Macros for moduli stacks/spaces
%
\def\QCohstack{\mathcal{QC}\!\mathit{oh}}
\def\Cohstack{\mathcal{C}\!\mathit{oh}}
\def\Spacesstack{\mathcal{S}\!\mathit{paces}}
\def\Quotfunctor{\mathrm{Quot}}
\def\Hilbfunctor{\mathrm{Hilb}}
\def\Curvesstack{\mathcal{C}\!\mathit{urves}}
\def\Polarizedstack{\mathcal{P}\!\mathit{olarized}}
\def\Complexesstack{\mathcal{C}\!\mathit{omplexes}}
% \Pic is the operator that assigns to X its picard group, usage \Pic(X)
% \Picardstack_{X/B} denotes the Picard stack of X over B
% \Picardfunctor_{X/B} denotes the Picard functor of X over B
\def\Pic{\mathop{\mathrm{Pic}}\nolimits}
\def\Picardstack{\mathcal{P}\!\mathit{ic}}
\def\Picardfunctor{\mathrm{Pic}}
\def\Deformationcategory{\mathcal{D}\!\mathit{ef}}
\newenvironment{reference}{}{}
\newenvironment{slogan}{}{}
\newenvironment{history}{}{}
\usepackage{cleveref}
\crefname{originalreference}{original reference}{original references}
\Crefname{originalreference}{Original reference}{Original references}
\makeatletter
\newlabel{divisors-definition-regular-section}{{01WY}{source}{}{sourcepointer-1}{}}
\newlabel{divisors-definition-regular-section@cref}{{[originalreference][0][]01WY}{[1][0][]source}{}{}{}}
\makeatother
\crefname{originallemma}{original lemma}{original lemmas}
\Crefname{originallemma}{Original lemma}{Original lemmas}
\makeatletter
\newlabel{lemma-proper-geometrically-reduced-global-sections}{{0BUG}{source}{}{sourcepointer-2}{}}
\newlabel{lemma-proper-geometrically-reduced-global-sections@cref}{{[originallemma][0][]0BUG}{[1][0][]source}{}{}{}}
\makeatother
\crefname{originallemma}{original lemma}{original lemmas}
\Crefname{originallemma}{Original lemma}{Original lemmas}
\makeatletter
\newlabel{lemma-vanishin-h0-negative}{{0FD7}{source}{}{sourcepointer-3}{}}
\newlabel{lemma-vanishin-h0-negative@cref}{{[originallemma][0][]0FD7}{[1][0][]source}{}{}{}}
\makeatother
\crefname{originallemma}{original lemma}{original lemmas}
\Crefname{originallemma}{Original lemma}{Original lemmas}
\makeatletter
\newlabel{lemma-vanishin-h1-negative}{{0FD8}{source}{}{sourcepointer-4}{}}
\newlabel{lemma-vanishin-h1-negative@cref}{{[originallemma][0][]0FD8}{[1][0][]source}{}{}{}}
\makeatother
\title{The zero scheme of a regular section of an ample invertible sheaf on a connected proper scheme of depth at least two is connected}
\author{The Stacks Project Authors}
\date{stacks-project:varieties.tex:local-snapshot-20261011; GNU FDL 1.2 or later}
\begin{document}
\maketitle
\noindent\textit{Editorial note.} All delivered mathematical text is the author's own, copied verbatim from the pinned Stacks Project chapter file varieties.tex; no statement or proof step was rewritten, reordered or removed.\par\medskip
\noindent\textit{Editorial note.} The counted span cites four cross-chapter labels (including a definition), each recorded in assembly.external\_references with its canonical tag as a navigation pointer.\par\medskip
\noindent\textit{Editorial note.} The proof cites no bibliography item inside the counted span.\par\medskip
\noindent\textit{Allowed standard prerequisites (editorial).}\begin{itemize}
\item Vanishing of H\textasciicircum{}0 and H\textasciicircum{}1 of negative twists of an ample invertible sheaf. Stacks Project, Varieties, Lemmas lemma-vanishin-h0-negative (tag 0FD7) and lemma-vanishin-h1-negative (tag 0FD8). The two vanishing statements contradicted by the disconnection of the zero scheme.
\item Regular sections and their zero schemes as effective Cartier divisors. Stacks Project, Divisors, Definition divisors-definition-regular-section (tag 0AYQ). Guarantees that the powers s\textasciicircum{}n are regular sections, which is what makes the comparison of zero schemes valid.
\end{itemize}\medskip
\ifdefined\hypertarget\hypertarget{sourcepointer-1}{}\fi
\noindent\textit{Source reference.} Original reference 01WY: sources/stacks-chow/tags.txt, source lines 2499--2499. Cross-chapter standard prerequisite cited by the author inside the counted span; it is used as a pointer and its own proof is not part of this item.\par\medskip
\ifdefined\hypertarget\hypertarget{sourcepointer-2}{}\fi
\noindent\textit{Source reference.} Original lemma 0BUG: sources/stacks-chow/tags.txt, source lines 14874--14874. Cross-chapter standard prerequisite cited by the author inside the counted span; it is used as a pointer and its own proof is not part of this item.\par\medskip
\ifdefined\hypertarget\hypertarget{sourcepointer-3}{}\fi
\noindent\textit{Source reference.} Original lemma 0FD7: sources/stacks-chow/tags.txt, source lines 19207--19207. Cross-chapter standard prerequisite cited by the author inside the counted span; it is used as a pointer and its own proof is not part of this item.\par\medskip
\ifdefined\hypertarget\hypertarget{sourcepointer-4}{}\fi
\noindent\textit{Source reference.} Original lemma 0FD8: sources/stacks-chow/tags.txt, source lines 19208--19208. Cross-chapter standard prerequisite cited by the author inside the counted span; it is used as a pointer and its own proof is not part of this item.\par\medskip
\begin{lemma}
\label{lemma-connectedness-ample-divisor}
Let $k$ be a field. Let $X$ be a proper scheme over $k$. Let $\mathcal{L}$
be an ample invertible $\mathcal{O}_X$-module. Let
$s \in \Gamma(X, \mathcal{L})$. Assume
\begin{enumerate}
\item $s$ is a regular section
(Divisors, Definition \ref{divisors-definition-regular-section}),
\item for every closed point $x \in X$ we have
$\text{depth}(\mathcal{O}_{X, x}) \geq 2$, and
\item $X$ is connected.
\end{enumerate}
Then the zero scheme $Z(s)$ of $s$ is connected.
\end{lemma}

\begin{proof}
Since $s$ is a regular section, so is
$s^n \in \Gamma(X, \mathcal{L}^{\otimes n})$ for all $n > 1$.
Moreover, the inclusion morphism $Z(s) \to Z(s^n)$ is a bijection
on underlying topological spaces. Hence if $Z(s)$ is disconnected,
so is $Z(s^n)$. Now consider the canonical short exact sequence
$$
0 \to \mathcal{L}^{\otimes -n} \xrightarrow{s^n}
\mathcal{O}_X \to \mathcal{O}_{Z(s^n)} \to 0
$$
Consider the $k$-algebra $R_n = \Gamma(X, \mathcal{O}_{Z(s^n)})$.
If $Z(s)$ is disconnected, i.e., $Z(s^n)$ is disconnected,
then either $R_n$ is zero in case $Z(s^n) = \emptyset$ or
$R_n$ contains a nontrivial idempotent in case $Z(s^n) = U \amalg V$ with
$U, V \subset Z(s^n)$ open and nonempty (the reader may wish to consult
Lemma \ref{lemma-proper-geometrically-reduced-global-sections}).
Thus the map $\Gamma(X, \mathcal{O}_X) \to R_n$ cannot be an isomorphism.
It follows that either $H^0(X, \mathcal{L}^{\otimes -n})$ or
$H^1(X, \mathcal{L}^{\otimes -n})$ is nonzero for infinitely
many positive $n$. This contradicts Lemma \ref{lemma-vanishin-h0-negative} or
\ref{lemma-vanishin-h1-negative}
and the proof is complete.
\end{proof}

\end{document}
